\graphicspath{ {../../images/} }
\usetikzlibrary{external}
\usepackage{changepage}

\title{ITC8280 Fundamentals of Cryptography}
\subtitle{Digital signatures \& PKI}
\date{\today}
\author{Taaniel Kraavi}
\institute%
{%
  \textit{School of Information Technologies}\\
  \textit{Tallinn University of Technology}
}

\begin{document}
\begin{frame}
  \titlepage
\end{frame}

\begin{frame}{Authenticity of data}
  Authenticity:
  \begin{itemize}[<+(1)->]
    \item Integrity protection: data modification can be detected
    \item Source authenticity: data sender can be identified
    \item Does not guarantee authorisation!
    \item Does not guarantee the correctness of the content!
  \end{itemize}

  \vspace*{1em}

  \pause
  MACs in the symmetric setting
  \begin{itemize}[<+(1)->]
    \item Key exchange problem
    \item \emph{Repudiation} problem
  \end{itemize}
\end{frame}

\begin{frame}{Digital signatures}
  Digital signatures:
  \begin{itemize}[<+(1)->]
    \item asymmetric / public-key scheme
    \begin{itemize}
      \item a public \emph{verification key} known to others
      \item a secret \emph{signing key} known to the signer
    \end{itemize}
    \item publicly verifiable: signatures are \emph{transferable}
    \item identification becomes possible
    \begin{itemize}
      \item keys are personal $\implies$ personal link
      \item identification is \emph{non-interactive} (NB! replay attacks)
      \item how informative is the identity? E.g. legal identification vs entity tracing
    \end{itemize}
    \item \emph{non-repudiation}: you cannot deny your signature
  \end{itemize}
\end{frame}

\begin{frame}{Digital signature scheme}
  A \emph{digital signature scheme} is a triple of PPT algorithms:
  \begin{itemize}[<+(1)->]
    \item $\GEN$ is a \emph{key generation algorithm}: $(\PK, \SK)\gets\GEN()$
    \item $\SIGN$ is a \emph{signing algorithm}: $\sigma\gets\SIGN_\SK(m)$
    \item $\VERIF$ is a deterministic \emph{verification algorithm}: $b\gets\VERIF_\PK(m, \sigma)$
  \end{itemize}

  \pause
  Notation and terms:
  \begin{itemize}[<+(1)->]
    \item $\SK$ is the private/secret key also called the \emph{signing key}
    \item $\PK$ is the public key also called the \emph{verification key}
    \item $b$ is the verification result (bit): $1$ meaning valid, $0$ meaning invalid
    \item $\sigma$ is a valid \emph{signature} on $m$ by $\SK$ if $b = 1$
  \end{itemize}

  \pause
  The signature scheme is \emph{functional} if $\VERIF_\PK(m, \SIGN_\SK(m)) = 1$.
\end{frame}

\begin{frame}{Digital signature scheme}
  \pause
  \begin{center}
    \begin{tikzpicture}
      \node (alice) at (0,0)  {\includegraphics[height=80px]{alice}};
      \node (bob)   at (8, 0) {\includegraphics[height=80px]{bob}};
      \node (eve)   at (4, 0) {\includegraphics[height=80px]{eve}};

      \draw[<-, thick] (eve.east) -- node[above] {$(m, \sigma$)} (bob.west);
      \draw[<-, thick] (alice.east) -- node[above] {$(\overline{m}, \overline{\sigma}$)} (eve.west);
      
      \node[anchor=east] (gen) at (11, 2.5) {$(\PK,\SK)\gets\GEN()$};
      \draw[<-,dashed,thick] (alice.north) |- (gen.west);

      \node (el) at (eve.north) [xshift=-5px] {};

      \path let \p1 = (el) in node (circ) at (\x1,2.5) {};
      \node at (circ.north) [yshift=5px] {$\PK$};

      \node[anchor=west]  at (8.5, 0.25)  {$m\gets\MMM$};
      \node[anchor=west]  at (8.5, -0.25) {$\sigma\gets\SIGN_\SK(m)$};
      \node[anchor=north] at (alice.south) {$\VERIF_\PK(\overline{m},\overline{\sigma})\iseq 1$};

      \draw[->,dashed,thick] (circ) -- (el);
      \node[anchor=center,draw,circle,fill=white,thick,minimum size = 2px, inner sep=0pt] at (circ) {};
    \end{tikzpicture}
  \end{center}

  \pause
  Mallory should not be able to \emph{forge} a valid signature.
\end{frame}

\begin{frame}{Security concepts}
  We consider \emph{chosen-message attacks} (CMA), which are stronger than KMA, KOA.

  \vspace*{1em}

  \pause
  Forgery types:
  \begin{itemize}[<+(1)->]
    \item \emph{Existential:}
    produce a valid signature for a new message
    \item \emph{Universal:}
    produce a valid signature for any message
  \end{itemize}

  \vspace*{1em}

  \pause
  A signature scheme should be \emph{EUF-CMA} secure.
  \begin{itemize}[<+(1)->]
    \item Strong existential unforgeability (sEUF) may sometimes be needed
    \item sEUF: Infeasible to forge a different signature for an already signed message
  \end{itemize}
\end{frame}

\begin{frame}{What to sign?}
  Signature schemes should handle arbitrary-length messages
  \begin{itemize}[<+(1)->]
    \item Messages may not fit within the mathematical bounds
    \item We need shorter message representatives
  \end{itemize}

  \vspace*{1em}

  \pause
  Hash-and-sign paradigm
  \begin{itemize}[<+(1)->]
    \item Cryptographic hashes are collision resistant
    \item Sign the hash of the message
    \item The hash is recomputed from the message during verification
  \end{itemize}
\end{frame}

\begin{frame}{Note about hashing}
  Do not blindly hash messages before signing
  \begin{itemize}[<+(1)->]
    \item Many signature schemes handle hashing internally
    \item \enquote{Pure} mode: input the full message; hashing is part of the algorithm
    \item \enquote{Pre-hash} mode: input an externally computed message digest
  \end{itemize}

  \vspace*{1em}

  \pause
  Prefer pure mode where possible
  \begin{itemize}[<+(1)->]
    \item Simpler: no need to keep track of additional parameters
    \item Security analysis covers the full message directly
  \end{itemize}
\end{frame}

\begin{frame}{RSA signatures}
  \pause
  PKCS \#1:
  \begin{itemize}[<+(1)->]
    \item \texttt{RSASSA-PKCS1-v1\_5} lacks tight security proofs, but not weak per se
    \item \texttt{RSASSA-PSS} (introduced in PKCS \#1 v2.1)
    \item \href{https://datatracker.ietf.org/doc/html/rfc8017}{\texttt{RFC 8017}}
  \end{itemize}

  \vfill

  \pause
  PSS -- Probabilistic Signature Scheme
  \begin{itemize}[<+(1)->]
    \item Should be used if RSA signatures are required
    \item Uses randomness (salt) to make signatures probabilistic
    \item Verification requires agreed parameters (hash, MGF, salt length)
  \end{itemize}
\end{frame}

\begin{frame}{RSASSA-PSS parameters}
  RSASSA-PSS: RSA Signature Scheme with Appendix with PSS padding
  \begin{itemize}[<+(1)->]
    \item An mask generation function (MGF, XOF) used internally like for RSAES-OAEP
    \item By default, the MGF is MGF1 (iterative hashing with a counter)
  \end{itemize}

  \vfill

  \pause
  Hashing in RSASSA-PSS:
  \begin{itemize}[<+(1)->]
    \item The message is internally hashed by the algorithm (SHA1 by default)
    \item The message is also hashed as part of MGF1 (SHA1 by default)
    \item You \emph{should} use the same hash function for the message and MGF1
    \item You \emph{should} use a SHA2 alternative instead of SHA1
  \end{itemize}
\end{frame}

\begin{frame}{Determinism vs. nondeterminisim}
  \pause
  Unlike for encryption, non-determinism is not crucial per se:
  \begin{itemize}[<+(1)->]
    \item no message to hide
  \end{itemize}

  \vspace*{0.5em}

  \pause
  Ephemeral randomness protects the key in many constructions:
  \begin{itemize}[<+(1)->]
    \item a leak might allow secret recovery
    \item reuse might allow secret recovery
  \end{itemize}

  \vspace*{0.5em}

  \pause
  Poor randomness is catastrophic in many schemes!
  \begin{itemize}[<+(1)->]
    \item (EC)DSA, Schnorr signatures
    \item Good randomness helps with uniqueness and secrecy
  \end{itemize}
\end{frame}

\begin{frame}{DSA \& ECDSA}
  \pause
  Digital Signature Algorithm (DSA):
  \begin{itemize}[<+(1)->]
    \item Hardness of the DLP
    \item Inspired by Schnorr signatures (historical patent reasons)
    \item Part of DSS versions 1--4, removed in version 5 (2023)
    \item \href{https://datatracker.ietf.org/doc/html/rfc6979}{\texttt{RFC 6979}}: deterministic DSA \& ECDSA
  \end{itemize}

  \pause
  Sony 2013 PS3 attack:
  \begin{itemize}
    \item ECDSA software signing-key recovered
    \item Caused by randomness reuse (Sony implementation issue)
  \end{itemize}
 
  \pause
  Do not use DSA! Do not default to \enquote{standard} ECDSA.
\end{frame}

\begin{frame}{EdDSA}
  Edwards-curve Digital Signature Algorithm (\href{https://datatracker.ietf.org/doc/html/rfc8032}{\texttt{RFC 8032}}):
  \begin{itemize}[<+(1)->]
    \item Not to be confused with DSA \& ECDSA: different algorithm entirely
    \item Designed by \alert{Bernstein} et al.
    \item Deterministic
    \item Based on twisted Edwards Curves (a model of ECs)
    \item High performance
    \item Designed with side-channel safety in mind
    \item Ed25519 (Curve25519 with SHA512) \& Ed448 (Curve448 with SHAKE256)
  \end{itemize}

  \pause
  Unless you have a good reason not to, use EdDSA!
\end{frame}

\begin{frame}{DSS (pt. 1)}
  NIST Digital Signature Standard (\href{https://csrc.nist.gov/pubs/fips/186-5/final}{FIPS 168-5}):
  \begin{itemize}[<+(1)->]
    \item No longer allows DSA for issuing signatures
    \begin{itemize}
      \item Prior signature verification still allowed
    \end{itemize}
    \item ECDSA is allowed
    \begin{itemize}
      \item Forbidden to use it for non-signing purposes (e.g. key establishment)
      \item Deterministic \& nondeterministic variants
    \end{itemize}
  \end{itemize}
\end{frame}

\begin{frame}{DSS (pt. 2)}
  NIST Digital Signature Standard (\href{https://csrc.nist.gov/pubs/fips/186-5/final}{FIPS 168-5}):
  \begin{itemize}[<+(1)->]
    \item EdDSA is allowed
    \begin{itemize}
      \item Both Ed25519 and Ed448 variants
      \item Variants differ a bit from the RFC
    \end{itemize}
    \item RSA signatures as defined in PKCS \#1 v2.2 are allowed
    \begin{itemize}
      \item Additional restrictions apply
      \item No textbook RSA!
    \end{itemize}
  \end{itemize}
\end{frame}

\begin{frame}{Long term signatures}
  How to archive digitally signed documents?
  \begin{itemize}[<+(1)->]
    \item In 20 years, how to check if a signature was valid today?
    \item Long-term validation (LTV)
    \item We will see formats and legality next week
  \end{itemize}

  \pause
  Re-signing over a document:
  \begin{itemize}[<+(1)->]
    \item Sign it with a new scheme before the old one is broken
    \item Timestamp the new signature for temporal order
    \item Extend your signature lifetime
  \end{itemize}

  \pause
  We will see timestamping also next week.
\end{frame}

\begin{frame}{Quantum-safe signature schemes}
  Quantum-safe schemes standardised by NIST:
  \begin{table}
  \small
  \begin{tabular}{@{}llll@{}}
    \toprule
    Name & Publication & Principle & When to pick?\\
    \midrule
    ML-DSA  & \href{https://csrc.nist.gov/pubs/fips/204/final}{FIPS 204} & Lattice & General purpose \\
    SLH-DSA & \href{https://csrc.nist.gov/pubs/fips/205/final}{FIPS 205} & Hash    & Conservative fallback \\
    FN-DSA  & FIPS 206 (draft)                                           & Lattice & Small signatures \\
    \bottomrule
  \end{tabular}
  \end{table}

  \vspace{0.5em}

  \pause
  PQ schemes are not yet widely used
  \begin{itemize}
    \item long keys and/or signatures
    \item slower than classical counterparts
  \end{itemize}
\end{frame}

\begin{frame}{Cool constructions}
  \pause
  Group signatures:
  \begin{itemize}
    \item Anonymously sign on behalf of a group
    \item The public cannot learn who this is
    \item The group \enquote{manager} can (unforgeably) trace and de-anonymise a signer
    \item Signatures are unlinkable (multiple signatures by the same issuer?)
  \end{itemize}

  \pause
  Ring signatures:
  \begin{itemize}
    \item Group signatures without a manager/trusted setup
    \item Traceable/linkable variants exist
  \end{itemize}

  \pause
  Blind signatures
  \begin{itemize}
    \item Sign some data without knowing the exact contents
  \end{itemize}
\end{frame}

\begin{frame}{Cool constructions}
  \pause
  Threshold signatures:
  \begin{itemize}
    \item Can only produce a valid signature in a quorum
  \end{itemize}

  \pause
  Multisignatures:
  \begin{itemize}
    \item Combine multiple signatures into a single succinct signature
    \item All signatories are evident from the multisignature
  \end{itemize}

  \pause
  Identity escrow:
  \begin{itemize}
    \item Party $A$ gives information to party $B$
    \item This information does not allow $B$ to identify $A$
    \item A 3\textsuperscript{rd} party $C$ can use this info to identify $A$
    \item $B$ receives a guarantee that $C$ can identify $A$
  \end{itemize}
\end{frame}

\begin{frame}{Trust models}
  How do we establish trust?
  \begin{itemize}[<+(1)->]
    \item Hierarchical trust (CAs)
    \item Web of trust (WoT)
    \item Trust on First Use (TOFU)
    \item Dual control model
    \item \dots
  \end{itemize}

  \pause
  Someone somewhere has to trust something sometime.
\end{frame}

\begin{frame}{Whomst'd've}
  How does Alice know that $\PK$ belongs to Bob?
  \begin{center}
    \begin{tikzpicture}
      \node (alice) at (0,0)  {\includegraphics[height=80px]{alice}};
      \node (bob)   at (8, 0) {\includegraphics[height=80px]{bob}};
      \node (eve)   at (4, 0) {\includegraphics[height=80px]{eve}};

      \draw[<-, thick] (eve.east) -- node[above] {$(m, \sigma$)} (bob.west);
      \draw[<-, thick] (alice.east) -- node[above] {$(\overline{m}, \overline{\sigma}$)} (eve.west);
      
      \node[anchor=east] (gen) at (11, 2.5) {$(\PK,\SK)\gets\GEN()$};
      \draw[<-,dashed,thick] (alice.north) |- (gen.west);

      \node (el) at (eve.north) [xshift=-5px] {};

      \path let \p1 = (el) in node (circ) at (\x1,2.5) {};
      \node at (circ.north) [yshift=5px] {$\PK$};

      \node[anchor=west]  at (8.5, 0.25)  {$m\gets\MMM$};
      \node[anchor=west]  at (8.5, -0.25) {$\sigma\gets\SIGN_\SK(m)$};
      \node[anchor=north] at (alice.south) {$\VERIF_\PK(\overline{m},\overline{\sigma})\iseq 1$};

      \draw[->,dashed,thick] (circ) -- (el);
      \node[anchor=center,draw,circle,fill=white,thick,minimum size = 2px, inner sep=0pt] at (circ) {};
    \end{tikzpicture}
  \end{center}
\end{frame}

\begin{frame}{Transitive trust}
  Consider the following:
  \begin{itemize}[<+(1)->]
    \item Carol is friends with both Alice and Bob
    \item Alice does not yet know Bob
    \item Friends know each other's public keys
  \end{itemize}

  \vspace*{1em}

  \pause
  Carol gives Bob her signature $\sigma_{C_B}$ on the message
  \begin{center}
    $C_B = $ \enquote{$\PK_B$ is Bob's public key}
  \end{center}
\end{frame}

\begin{frame}{Transitive trust}
  \begin{figure}
    \centering
    \begin{tikzpicture}
      \node (alice) at (0,0)  {\includegraphics[height=80px]{alice}};
      \node (bob)   at (8, 0) {\includegraphics[height=80px]{bob}};

      \draw[<-, thick] (alice.east) -- node[above] {$(m, \sigma),\ \bigl(C_B, \sigma_{C_B}\bigr)$} (bob.west);

      \node[anchor=west]  at (8.5, 0.25)  {$m\gets\MMM$};
      \node[anchor=west]  at (8.5, -0.25) {$\sigma\gets\SIGN_{\SK_B}(m)$};
      \node[anchor=north] at (alice.south) {\small$\VERIF_{\PK_C}(\PK_B,\sigma_{\PK_B})\iseq 1$};
      \node[anchor=north] at (alice.south) [yshift=-1.5em] {\small$\VERIF_{\PK_B}(m,\sigma)\iseq 1$};

      \node (alice-left) at (-2, 2) {};
      \draw[->,dashed,thick] (alice-left) -- node[above] {$\PK_C$} (alice-left -| alice.north) -- (alice.north);

      \node (bob-right) at (10, 2) {};
      \draw[->,dashed,thick] (bob-right) -- node[above] {$C_B,\ \sigma_{C_B}$} (bob-right -| bob.north) -- (bob.north);
    \end{tikzpicture}
  \end{figure}

  \pause
  $C_B$ can be viewed as a \emph{certificate} by Carol linking Bob's public key to him.
\end{frame}

\begin{frame}{Digital certificates}
  A \emph{certificate} binds keys to an entity's attributes (e.g. name, email).

  \vspace*{1em}

  \pause
  Let there be two parties: a \emph{subject} and an \emph{issuer}.
  \begin{enumerate}[<+(1)->]
    \item The subject sends their public key and identity attributes to the issuer
    \item The issuer verifies the info, bundles it together, and signs the bundle
    \begin{itemize}
      \item The issuer confirms the attributes somehow
      \item The issuer verifies secret key ownership of the entity
    \end{itemize}
  \end{enumerate}

  \vspace*{1em}

  \pause
  The signed bundle is the certificate.
\end{frame}

\begin{frame}{Role of the issuer}
  The issuer acts like a notary.
  \begin{itemize}[<+(1)->]
    \item You \enquote{simply} trust the notary
    \item Trust the notary $\implies$ trust their issued certificate
  \end{itemize}

  \vspace*{1em}

  \pause
  Verification requirements vary
  \begin{itemize}[<+(1)->]
    \item e.g.\ internal company use vs cross-border trust
    \item e.g.\ bank transaction vs git commit
  \end{itemize}
\end{frame}

\begin{frame}{Public Key Infrastructure (PKI)}
  A public key infrastructure consists of a:
  \begin{itemize}[<+(1)->]
    \item certificate authority (CA): stores, issues \& signs certificates
    \item registration authority (RA): verifies the identity of certificate requesters
    \item central directory: indexing of keys
  \end{itemize}

  \vspace*{1em}

  \pause
  Entities need not all be separate, e.g. CA \& RA could be one.
  \begin{itemize}
    \item The term \emph{trusted third party} (TTP) often used
  \end{itemize}
\end{frame}

\begin{frame}{PKI (extended)}
  \pause
  There is also usually a:
  \begin{itemize}[<+(1)->]
    \item certificate policy: defines parties' roles and procedures
    \item certificate management system
  \end{itemize}

  \pause
  These are less critical for understanding how the PKI works.
\end{frame}

\begin{frame}{Root CAs}
  \emph{Root CAs} are trust anchors:
  \begin{itemize}[<+(1)->]
    \item \emph{Root certificates}: self-signed certificates issued by a root CA
    \item Root certificates must be obtained from trusted sources
    \par
    \begin{itemize}
      \item Root certificates are often bundled with software, e.g.\ OS and browsers
      \item Users may manually add root certificates to their certificate store
    \end{itemize}
  \end{itemize}

  \vspace*{1em}

  \pause
  Root private keys must be well protected and are often kept offline.
  \begin{itemize}[<+(1)->]
    \item Not used to issue \emph{end-entity certificates} directly
    \item Key change/rotation $\implies$ new trust anchor must be distributed
  \end{itemize}
\end{frame}

\begin{frame}{Intermediate CAs}
    Trust is extended through \emph{intermediate CAs}:
    \begin{itemize}[<+(1)->]
      \item Certified by a root CA or another intermediate CA
      \item Can issue certificates to other intermediates or end entities
      \item Handle day-to-day issuance to better handle risk
    \end{itemize}

    \vspace*{1em}

    \pause
    If an intermediate key is compromised:
    \begin{itemize}[<+(1)->]
      \item \enquote{Only} its subtree of certificates is affected
      \item The issuer can revoke and replace the intermediate
    \end{itemize}
\end{frame}

\begin{frame}{End-Entities}
  \emph{End-entity certificates} terminate the trust chain:
  \begin{itemize}[<+(1)->]
    \item End-entities: a user, device, service, \dots
    \par
    \begin{itemize}
      \item e.g.\ website, digital signature author, software packager
    \end{itemize}
    \item Also called \emph{leaf} certificates
    \item Certified by an intermediate CA
    \item Cannot issue further certificates
  \end{itemize}

  \vspace*{1em}

  \pause
  End-entity key compromise does not affect the rest of the chain.
\end{frame}

\begin{frame}{Chain of Trust}
  Certificate chain:
  \begin{center}
    root cert $\to$ intermediate cert(s) $\to$ end-entity cert
  \end{center}

  \vspace*{1em}

  \pause
  To verify a certificate, verify the chain up to a trusted root
  \begin{itemize}[<+(1)->]
    \item Both signatures and certificate properties must be checked
    \item Root certificates must be obtained from trusted sources
    \item If a private key leaks, all certificates chaining to it become untrustworthy
  \end{itemize}
\end{frame}

\begin{frame}{X.509 certificates}
  \pause
  Bind a public key to an identity using a digital signature.
  \begin{itemize}[<+(1)->]
    \item \href{https://www.itu.int/rec/T-REC-X.509}{ITU-T X.509}
    \item \enquote{For the Internet}: \href{https://datatracker.ietf.org/doc/html/rfc5280}{RFC 5280}
    \item \href{https://cabforum.org}{CA/Browser Forum certificate baselines}
    \item Self-signed or signed by a CA
    \item ASN.1 structures
  \end{itemize}
\end{frame}

\begin{frame}{X.509 v3 structure}
  \begin{itemize}[<+(1)->]
    \item Certificate
    \begin{itemize}
      \item Version \& serial number
      \item Issuer
      \item Subject
      \item Validity period: not before \& not after
      \item Subject Public Key Info: algorithm \& public key
      \item Signature Algorithm ID
      \item Issuer and/or Subject Unique Identifier (optional)
      \item Extensions (optional)
    \end{itemize}
    \item Signature Algorithm
    \item Signature Value (of the certificate by the issuer)
  \end{itemize}

  \pause
  The above structure is simplified and reordered.
\end{frame}

\begin{frame}{Certificate extensions}
  Certificate extensions:
  \begin{itemize}[<+(1)->]
    \item Uniquely identified by an Object Identifier (OID)
    \item Labelled \enquote{critical} or \enquote{non-critical}
    \par
    \begin{itemize}
      \item Must reject certificate if unrecognised critical extension
      \item May ignore an unrecognised non-critical extension
      \item Must process extensions if recognised
    \end{itemize}
    \item Contain a set of values
  \end{itemize}

  \pause
  Restrict usage for specific purposes
  \begin{itemize}[<+(1)->]
    \item All restrictions must be satisfied
    \item Downgrade attacks might be possible
  \end{itemize}
\end{frame}

\begin{frame}{Certificate extensions}
  Some examples:
  \begin{itemize}[<+(1)->]
    \item Basic constraints
    \begin{itemize}
      \item Whether the subject is a CA
      \item Maximal certification path depth
    \end{itemize}
    \item Key usage
    \begin{itemize}
      \item Allowed cryptographic operations
      \item E.g. certificate signing only
    \end{itemize}
    \item Extended Key Usage
    \begin{itemize}
      \item Fine-tune key usage, e.g. only allowed for email
      \item Must not conflict with key usage
    \end{itemize}
    \item \dots
  \end{itemize}
\end{frame}

\begin{frame}{Web of Trust}
  An alternative to the PKI chain of trust:
  \begin{itemize}[<+(1)->]
    \item Most widely used with PGP
    \item Not centralised unlike CA-based systems (PKIs in general)
    \begin{itemize}
      \item LinkedIn for public keys
      \item Upload your keys to key-servers on the web (no deletion possible usually)
      \item Maintain a local database of others' public keys (the \emph{keyring})
    \end{itemize}
    \item You \emph{vet} a user by signing their keys
    \begin{itemize}
      \item Others who trust you can then trust the keys you have signed
      \item This is done by configuring key trust levels
    \end{itemize}
  \end{itemize}

  \pause
  It's a mess, especially the key management.
\end{frame}

\end{document}
